\section{Билет 68. Термодинамическое определение энтропии. Вычисление приращения энтропии для изопроцессов идеального газа: изотермического и изохорического.}

\begin{definition}
\textbf{Энтропия} $S$ --- функция состояния, характеризующая меру необратимости рассеяния энергии в термодинамической системе. Для обратимого процесса:
\begin{equation}
    dS = \frac{\delta Q_{\text{обр}}}{T}, \quad \Delta S = \int_1^2 \frac{\delta Q_{\text{обр}}}{T}. \label{eq:entropy_def_68}
\end{equation}
Единица измерения энтропии в СИ --- джоуль на кельвин (Дж/К).
\end{definition}

\begin{theorem}[Приращение энтропии в изотермическом процессе]
При $T = \text{const}$ для идеального газа $\Delta U = 0$, поэтому $\delta Q = p\,dV$. Используя уравнение состояния $p = \nu RT/V$:
\begin{equation}
    \Delta S_{\text{изот}} = \int_1^2 \frac{p\,dV}{T} = \nu R \int_{V_1}^{V_2} \frac{dV}{V} = \nu R \ln\frac{V_2}{V_1}. \label{eq:delta_S_iso_68}
\end{equation}
Через давления (из $p_1 V_1 = p_2 V_2$):
\begin{equation}
    \Delta S_{\text{изот}} = \nu R \ln\frac{p_1}{p_2}.
\end{equation}
\end{theorem}

\begin{theorem}[Приращение энтропии в изохорическом процессе]
В изохорическом процессе объём постоянен: $V = \text{const}$, поэтому работа $A = 0$ и $\delta Q_V = dU = \nu C_V \, dT$, где $C_V$ --- молярная теплоёмкость при постоянном объёме.
Подставляя в определение энтропии:
\begin{equation}
    dS = \frac{\delta Q_V}{T} = \nu C_V \frac{dT}{T}.
\end{equation}
Интегрируя от $T_1$ до $T_2$:
\begin{equation}
    \Delta S_{\text{изох}} = \nu C_V \int_{T_1}^{T_2} \frac{dT}{T} = \nu C_V \ln\frac{T_2}{T_1}. \label{eq:delta_S_isoch_68}
\end{equation}
Используя закон Шарля $p/T = \text{const}$, энтропию можно выразить через давления:
\begin{equation}
    \Delta S_{\text{изох}} = \nu C_V \ln\frac{p_2}{p_1}.
\end{equation}
\end{theorem}

\begin{property}[Сравнение формул для изопроцессов]
\begin{center}
\begin{tabular}{|l|l|l|}
\hline
Процесс & Условие & $\Delta S$ \\
\hline
Изотермический & $T=\text{const}$ & $\nu R \ln\dfrac{V_2}{V_1} = \nu R \ln\dfrac{p_1}{p_2}$ \\
\hline
Изобарический & $p=\text{const}$ & $\nu C_p \ln\dfrac{T_2}{T_1} = \nu C_p \ln\dfrac{V_2}{V_1}$ \\
\hline
Изохорический & $V=\text{const}$ & $\nu C_V \ln\dfrac{T_2}{T_1} = \nu C_V \ln\dfrac{p_2}{p_1}$ \\
\hline
\end{tabular}
\end{center}
\end{property}

\begin{remark}
\textbf{Важные замечания:}
\begin{enumerate}
    \item Все формулы справедливы только для \textbf{обратимых} процессов. Для необратимых процессов энтропия вычисляется по тому же интегралу, но вдоль воображаемого обратимого пути, соединяющего те же начальное и конечное состояния.
    \item Для идеального газа $C_p = C_V + R$ и $\gamma = C_p/C_V$, поэтому $\Delta S_{\text{изоб}} > \Delta S_{\text{изох}}$ при одинаковом изменении температуры, так как $C_p > C_V$.
    \item В адиабатическом обратимом процессе $\delta Q = 0$, поэтому $\Delta S = 0$ --- энтропия сохраняется. Такие процессы называются \textbf{изоэнтропийными}.
\end{enumerate}
\end{remark}
